% Generated draft: requires mathematical and source review.
\documentclass[11pt]{article}
\usepackage[margin=25mm]{geometry}
\usepackage{fontspec}
\setmainfont{Latin Modern Roman}
\usepackage{amsmath,amssymb,mathtools,unicode-math}
\setmathfont{Latin Modern Math}
\usepackage{enumitem,needspace,xurl,hyperref}
\setlength{\parindent}{0pt}
\setlength{\parskip}{5pt}
\setlength{\emergencystretch}{3em}
\newcommand{\N}{\mathbb{N}}
\newcommand{\Z}{\mathbb{Z}}
\newcommand{\Q}{\mathbb{Q}}
\newcommand{\R}{\mathbb{R}}
\newcommand{\C}{\mathbb{C}}
\newcommand{\F}{\mathbb{F}}
\newcommand{\A}{\mathbb{A}}
\newcommand{\PP}{\mathbb{P}}
\newcommand{\E}{\mathbb{E}}
\newcommand{\eps}{\varepsilon}
\newcommand{\dd}{\,\mathrm{d}}
\begin{document}
\section*{UnsolvedMath ID: 30005083}
\textbf{Source problem code:} catalogue entry
\subsection{Definitions and mathematical statement}
The supplied registry record concerns the nondegeneracy assertion in the spectral edges conjecture. Its reconstructed statement is: for a sufficiently general periodic discrete Schrödinger operator \(L\), every extremum of the dispersion function \(\lambda\) on the real Bloch variety is nondegenerate. The operator class and the meaning of ``sufficiently general'' are not specified in the supplied evidence, so this does not yet determine a fully quantified conjecture.

The following gives a conditional framework for interpreting the assertion. Suppose that the periodic operator under consideration has a finite-dimensional Hermitian Floquet matrix \(L(k)\), smoothly parametrized by quasimomentum \(k\) in a \(d\)-dimensional torus \(\mathbb T^d\). Here \(d\geq 1\), the matrix size, and the permitted operators must be fixed by the intended problem. In this framework define the real Bloch variety and energy projection by
\[
 B_{\mathbb R}(L)=\{(k,E)\in\mathbb T^d\times\mathbb R:\det(L(k)-EI)=0\},
 \qquad \lambda(k,E)=E.
\]
At a point \(p\) where \(B_{\mathbb R}(L)\) is a smooth manifold, an extremum means, under the local interpretation, that \(\lambda(p)\) is a local maximum or local minimum of the restricted function. Such an extremum is nondegenerate if the Hessian of \(\lambda|_{B_{\mathbb R}(L)}\) at \(p\), computed in local coordinates, is nonsingular. This condition is independent of the coordinate choice because the differential vanishes at an extremum.

In particular, near a simple Floquet eigenvalue the Bloch variety is locally the graph of a smooth eigenvalue branch \(E_j(k)\). At a local extremum \(k_0\), the condition becomes
\[
 \nabla_k E_j(k_0)=0,
 \qquad
 \det\!\left(\frac{\partial^2 E_j}{\partial k_a\partial k_b}(k_0)\right)_{a,b=1}^{d}\neq 0.
\]
This is the full Hessian in all quasimomentum directions.

A precise genericity formulation would require an explicitly specified parameter space \(\mathcal P\) of periodic operators and a specified meaning of a general subset \(\mathcal G\subseteq\mathcal P\). Its intended conclusion would then quantify over every \(L\in\mathcal G\) and every relevant extremum of \(\lambda\). Neither \(\mathcal P\), the required size or type of \(\mathcal G\), nor the treatment of singular points and eigenvalue crossings is supplied. The smooth-point definition above therefore explains nondegeneracy without resolving those missing parts of the assertion. It also does not establish whether the intended extrema are all local extrema or only those associated with spectral edges.

The supplied original excerpt merely identifies understanding critical points of \(\lambda\) as a first step toward the spectral edges conjecture and includes a figure caption with edge parameters \((6,3,2)\) and \((1,1,1)\). It does not itself state the reconstructed nondegeneracy assertion or specify the underlying graph.
\subsection{Short English statement}
The reconstructed conjecture asks whether every relevant maximum and minimum of the dispersion relation has a nonsingular Hessian for sufficiently general periodic discrete Schrödinger operators. The permitted operators, the meaning of generality, and the treatment of crossings and singular points require clarification. This is a statement-only draft requiring human mathematical review.
\subsection{Sources}
\begin{itemize}
\item Supplied local UnsolvedMath registry record, permanent ID \(30005083\), title ``Nondegeneracy in the Spectral Edges Conjecture,'' source problem code \texttt{catalogue entry}, in the Oberwolfach Reports Open Problems collection. No individual proposer is supplied. The record gives proposed year \(2022\).
\item Supplied report citation: \emph{Toric Geometry (2023)}, Oberwolfach Reports, DOI \texttt{10.4171/owr/2022/17}, pp.\ 861--927. \url{https://doi.org/10.4171/owr/2022/17}. The supplied title uses \(2023\), while the proposed year and DOI contain \(2022\); these details are retained as supplied.
\item The registry labels the statement \texttt{corrected\_verified} and says it was reconstructed and checked against the official Oberwolfach report. This is a supplied verification claim, not an independent verification performed for this draft. The supplied original excerpt is a discussion fragment and figure caption.
\item Supplied literature reference: Faust--Sottile, \emph{Journal of Spectral Theory}. \url{https://doi.org/10.4171/JST/503}.
\item Supplied literature reference: Faust--Robinson--Sottile. \url{https://arxiv.org/abs/2512.01226}.
\item Literature assessment dated exactly \texttt{2026-08-22}: the supplied curation labels the problem partially solved and reports that critical-point bounds and effective criteria are proved, that the spectral-edges conjecture is established for broad families of periodic graphs, and that the general conjecture remains unresolved. This is supplied dated curation derived from the OWR record, not an independently verified assessment; its time-sensitive status and cited claims require independent checking.
\end{itemize}
\end{document}
